\documentclass[journal]{IEEEtran}
\usepackage{amsmath}
\usepackage{amssymb}
\usepackage{bm}
\usepackage{amsfonts}
\usepackage{upgreek}
\usepackage[sorting=none]{biblatex} %Imports biblatex package
% \addbibresource{regularized-maximum-likelihood-estimation.bib}
\addbibresource{references.bib}
\begin{document}
% Linebreaks \\ can be used within to get better formatting as desired.
\title{Overview of AFN algorithms}
\author{Kieran~Choi-Slattery}
% \usepackage{mhchem}
% The paper headers
\markboth{SSDL Planetary eXploration Lab - June 14, 2024}%
{Shell \MakeLowercase{\textit{et al.}}: 
%Rewrite the title here if you have many pages and want it to appear on other ones in the corner
}

% make the title area
\maketitle
\newcommand{\erf}{\mathrm{erf}}
\newcommand{\Cp}{C_{\scriptscriptstyle{P}}}
\newcommand{\Cs}{C_{\scriptscriptstyle{S}}}
\newcommand{\sigmap}{\sigma_{\scriptscriptstyle{P}}}
\newcommand{\sigmas}{\sigma_{\scriptscriptstyle{S}}}
\section{Nomenclature}
\begin{tabular}{l l}
    Symbol & Definition \\
    \hline
    $P, p$ & Parallel power collected by AFN\\
    $S, s$ & Perpendicular power collected by AFN\\
    $C_{\scriptstyle{P}}$ & Parallel scattering cross-section\\
    $C_{\scriptstyle{S}}$ & Perpendicular scattering cross-section\\
    $D, \delta$ & Particle diameter\\
    $m$ & Complex refractive index of particle\\
    $N, n$ & Real part of particle refractive index\\
    $k$ & Imaginary refractive index, $m=n-ik$\\
    $B$ & Incident irradiance (power/area) on particle\\
\end{tabular}
\section{Parallel and perpendicular cross-sections}
The cross-sections have units of area and represent the quantity that, when multiplied by the incident irradiance $B$, return the power incident on the AFN optics due to scattering from the particle:
\begin{equation}
\begin{bmatrix}
    P\\
    S
\end{bmatrix}=
B
\begin{bmatrix}
    C_{\scriptstyle{P}}\\
    C_{\scriptstyle{S}}
\end{bmatrix}
\end{equation}
They are calculated using the equations:
\begin{align}
    C_{\scriptstyle{P}}&=A\int_{\Omega_{col}}|S_2(\theta, x, m)|^2\ d\Omega \label{eq:sigmap}\\
    C_{\scriptstyle{S}}&=A\int_{\Omega_{col}}|S_1(\theta, x, m)|^2\ d\Omega. \label{eq:sigmas}
\end{align}
where $A$ is the particle's geometric cross-section,
\begin{equation}
    A=\frac{\pi}{4}D^2.
\end{equation}
$S_1$ and $S_2$ are calculated using the methods described in, e.g., \cite{bohren_absorption_2004} or \cite{prahl_miepython_2023}, but normalized such that, when integrated around all angles, they are equal to the particle's total scattering efficiency $Q_{sca}$, which represents the ratio of total scattered power to power incident on the particle:
\begin{equation}
    \int_{4\pi}|S_1(\theta, x, m)|^2+|S_2(\theta, x, m)|^2\ d\Omega=Q_{sca}.
\end{equation}
$x$ represents the size parameter,
\begin{equation}
    x=\frac{\pi D}{\lambda},
\end{equation}
with $\lambda$ the excitation wavelength (outside of the particle).
\par The integrations in equations \ref{eq:sigmap} and \ref{eq:sigmas} are performed by integrating in spherical coordinates and making the substitution that $d\Omega$ represents a differential solid angle, $d\Omega=\sin{\theta}d\phi d\theta$. 
\par For an arbitrary collector with an angular profile defined by a window $\nu(\theta)=\phi_{max}(\theta)-\phi_{min}(\theta)$,
\begin{align}
    C_{\scriptstyle{P}}&=A\int_{\theta_{min}}^{\theta_{max}}|S_2(\theta, x, m)|^2\nu(\theta)\sin{\theta}d\theta \label{eq:sigmapint}\\
    C_{\scriptstyle{S}}&=A\int_{\theta_{min}}^{\theta_{max}}|S_1(\theta, x, m)|^2\nu(\theta)\sin{\theta}d\theta \label{eq:sigmasint}
\end{align}
% \par For an off-axis circular collector with half-aperture angle $\alpha$ centered on a point with an angle $\theta_c$ made from the forward direction of the beam to the line between the particle and the collector, 
\par For a collector with the geometry of the AFN, the window function $\nu(\theta)$ is given by \cite{hodkinson_response_1965} as:
\begin{equation}\label{weightfn}
    \nu(\theta)=2\arccos\left(\frac{\cos \alpha -\cos \theta_c\cos \theta }{\sin \theta_c \sin \theta }\right)
\end{equation}
with $\alpha$ the angular half-aperture of the circular collector and $\theta_c$ the polar angle of the center of the collector.
\par For computation, equations \ref{eq:sigmapint} and \ref{eq:sigmasint} must be discretized:
\begin{align}
    C_P&=\frac{\pi D^2}{4}\sum_{i}|S_2(\theta_i, x, m)|^2 \nu(\theta_i)\sin{\theta_i}\ \Delta\theta\\
    C_S&=\frac{\pi D^2}{4}\sum_{i}|S_1(\theta_i, x, m)|^2 \nu(\theta_i)\sin{\theta_i}\ \Delta\theta
\end{align}

\section{Ensembles of particles}
\par Let $W_{\scriptscriptstyle{P}}$ and $W_{\scriptscriptstyle{S}}$ be random variables representing additive noise to the signal. $W_{\scriptscriptstyle{P}}$, $W_{\scriptscriptstyle{S}}$, $D$, $N$, and $B$ are all random continuous variables. Let also $\mathbf{U}=\begin{bmatrix}
    P & S
\end{bmatrix}^\intercal$ be the measurement vector and $\mathbf{C}(\delta, n): \mathbb{R}^2\mapsto \mathbb{R}^2$ be the function that maps a diameter and refractive index to a pair of cross-sections, so that
\begin{equation}
    \mathbf{U}=B\mathbf{C}(D, N)+\begin{bmatrix}W_{\scriptscriptstyle{P}}\\W_{\scriptscriptstyle{S}}\end{bmatrix}
\end{equation}
When a particle is sampled from the ensemble, it is picked randomly from the joint probability density function $f(w_{\scriptscriptstyle{P}}, w_{\scriptscriptstyle{S}},\delta, n, b)$. The probability that $\mathbf{U}$ falls within a specific set of values $\mathcal{U}$ is given by
\begin{multline}\label{eq:ubinprob}
    \mathrm{Pr}[\mathbf{U}\in \mathcal{U}]=\\
    \int\!\!\underset{\mathrm{\times5}}{\cdots}\!\!\int1_{\mathcal{U}}\left(b\mathbf{C}(\delta, n)+\begin{bmatrix}w_{\scriptscriptstyle{P}}\\w_{\scriptscriptstyle{S}}\end{bmatrix}\right)\\
    f(w_{\scriptscriptstyle{P}}, w_{\scriptscriptstyle{S}}, \delta, n, b)\,dw_{\scriptscriptstyle{P}}\,dw_{\scriptscriptstyle{S}}\,db\,dn\,d\delta
\end{multline}
where the integrals are across all possible values of the integration variables and $1_{\mathcal{U}}(\mathbf{u})$ is the indicator function of $\mathcal{U}$:
\begin{equation}
    1_{\mathcal{U}}(\mathbf{u})=\begin{cases}
        1 & \textrm{if } \mathbf{u} \in \mathcal{U}\\
        0 & \textrm{otherwise}
    \end{cases}
\end{equation} 
\par Since $W_{\scriptscriptstyle{P}}$, $W_{\scriptscriptstyle{S}}$, and $B$ are all independent variables, 
\begin{multline}
    f(w_{\scriptscriptstyle{P}}, w_{\scriptscriptstyle{S}}, \delta, n, b)=\\
    f_{\scriptscriptstyle{D, N}}(\delta, n)\cdot f_{\scriptscriptstyle{B}}(b)\cdot f_{\scriptscriptstyle{W_{\scriptscriptstyle{P}}}}(w_{\scriptscriptstyle{P}})\cdot
    f_{\scriptscriptstyle{W_{\scriptscriptstyle{S}}}}(w_{\scriptscriptstyle{S}}).
\end{multline}
Substituting into equation \ref{eq:ubinprob},
\begin{multline}\label{eq:ubinprob2}
    \mathrm{Pr}[\mathbf{U}\in \mathcal{U}]=
    \iint \mathrm{Pr}[\mathbf{U}\in\mathcal{U}|\delta, n]\,f_{\scriptscriptstyle{D, N}}(\delta, n)\,dn\,d\delta
\end{multline}
where $\mathrm{Pr}[\mathbf{U}\in\mathcal{U}|\delta, n]$ represents that particle with diameter $\delta$ and refractive index $n$ produces a measurement within $\mathcal{U}$, given by
\begin{equation}
    \mathrm{Pr}[\mathbf{U}\in\mathcal{U}|\delta, n]=\int \mathrm{Pr}[\mathbf{U}\in\mathcal{U}|\delta, n, b]f_{\scriptscriptstyle{B}}(b) \, db
\end{equation}
where $\mathrm{Pr}[\mathbf{U}\in\mathcal{U}|\delta, n, b]$ is the probability that the aforementioned particle, having fallen within a part of the beam with intensity $b$, produces a measurement within $\mathcal{U}$.
\par For normally-distributed $W_{\scriptscriptstyle{P}}$, $W_{\scriptscriptstyle{S}}$, with standard deviations $\sigmap$, $\sigmas$ and rectangular $\mathcal{U}$ defined by lower-bounds ($p_0$, $s_0$) and upper-bounds ($p_1$, $s_1$), $\mathrm{Pr}[\mathbf{U}\in\mathcal{U}|\delta, n, b]$ is given analytically as
\begin{multline}
    \mathrm{Pr}[\mathbf{U}\in\mathcal{U}|\delta, n, b]=\\ \frac{1}{4}\left(\erf\left(\frac{p_1-b\Cp}{\sigmap\sqrt{2}}\right)-\erf\left(\frac{p_0-b\Cp}{\sigmap\sqrt{2}}\right)\right)\\
    \cdot\left(\erf\left(\frac{s_1-b\Cs}{\sigmas\sqrt{2}}\right)-\erf\left(\frac{s_0-b\Cs}{\sigmas\sqrt{2}}\right)\right)
\end{multline}

The transformation from the space of probability density functions to the space of measured histograms is thus a linear operator which acts on the function $f_{\scriptscriptstyle{D, N}}(\delta, n)$. We denote this linear operator as $\bm{R}$, so that
\begin{equation}
    \bm{R}f_{\scriptscriptstyle{D, N}}=\begin{bmatrix}
        \mathrm{Pr}[\mathbf{U}\in\mathcal{U}_1]\\
        \mathrm{Pr}[\mathbf{U}\in\mathcal{U}_2]\\
        \vdots
    \end{bmatrix}
\end{equation}
with $\mathcal{U}_1$, $\mathcal{U}_2$, etc. the bins chosen for the 2D histogram.

\section{Numerical representation of $f_{\scriptscriptstyle{D, N}}$}
In order to invert the measured data to find $f_{\scriptscriptstyle{D, N}}(\delta, n)$, it is decomposed into a linear combination of basis functions ${\psi(\delta, n)}$ with coefficients $\{a\}$:
\begin{equation}
    f_{\scriptscriptstyle{D, N}}(\delta, n)\approx\sum_{j}a_j\psi_{j}(\delta, n)
\end{equation}
This transformation can be represented as a linear operator $\bm{\Psi}$:
\begin{equation}
    f_{\scriptscriptstyle{D, N}}\approx\bm{\Psi}\mathbf{a}\textrm{, }\mathbf{a}=\begin{bmatrix}a_1\\a_2\\\vdots\end{bmatrix}
\end{equation}
Because $\mathbf{C}(\delta, n)$ is extremely sensitive to both $\delta$ and $n$, repeatedly performing the integration in equation \ref{eq:ubinprob2} to a reasonable precision is infeasible. However, due to the linearity of equation \ref{eq:ubinprob2}, 
\begin{equation}
    \mathrm{Pr}[\mathbf{U}\in \mathcal{U}_i]\approx\sum_j a_j\cdot\mathrm{Pr}[\mathbf{U}\in \mathcal{U}_i|\psi_j]
\end{equation}
where
\begin{equation}\label{eq:rpsi}
    \mathrm{Pr}[\mathbf{U}\in \mathcal{U}_i|\psi_j]=\iint \mathrm{Pr}[\mathbf{U}\in\mathcal{U}_i|\delta, n]\psi_j(\delta, n)\,dn\,d\delta.
\end{equation}
Each integral can be evaluated individually, before the inversion, and then the operation of evaluating a guess for $f_{\scriptscriptstyle{D, N}}$ can be performed as a matrix operation:
\begin{equation}
    \bm{R}\bm{\Psi}\mathbf{a}\approx\begin{bmatrix}
        \mathrm{Pr}[\mathbf{U}\in\mathcal{U}_1]\\
        \mathrm{Pr}[\mathbf{U}\in\mathcal{U}_2]\\
        \vdots
    \end{bmatrix}
\end{equation}
If non-negative functions are chosen as the basis, then by constraining $\mathbf{a}$ to be non-negative, all possible values of $\bm{\Psi}\mathbf{a}$ will be as well.
\section{Inverse problem}
In the inverse problem, data is collected and binned. The number of particles that fell into each bin are recorded as
\begin{equation}
\mathbf{h}=\begin{bmatrix}h_1\\h_2\\\vdots\end{bmatrix},
\end{equation}
where $h_i$ represents the number of particles that fell within the $i$'th bin. At a high level, the inverse algorithm can be stated as
\begin{multline}
    \hat{f}_{\scriptscriptstyle{D, N}}=\bm{\Psi}\mathbf{\hat{a}},\ {\mathbf{\hat{a}}}=\underset{\mathbf{a}}{\operatorname{argmin}}\ \Delta(\eta\bm{R}\bm{\Psi}\mathbf{a}, \mathbf{h})+\epsilon L(\bm{\Psi}\mathbf{a})\\
    \textrm{such that} \begin{cases}
        \operatorname{min}\mathbf{a}\geq0\\
        \iint\bm{\Psi}\mathbf{a}\, dn\,d\delta=1
    \end{cases}
\end{multline}
where 
\begin{itemize}
    \item $\eta$ is the number of particles observed.
    \item $\Delta(\bm{R}\bm{\Psi}\mathbf{a}, \mathbf{h})$ is some statistical divergence between the probabilities produced by $\bm{R}\bm{\Psi}\mathbf{a}$ (the model) and measured by $\mathbf{h}$ (the data).
    \item $L(\bm{\Psi}\mathbf{a})$ is the regularizer loss function, which quantifies how non-smooth the argument is.
    \item $\epsilon$ is the regularization parameter, which determines how much the loss function should be weighted.
\end{itemize}
\subsection{Divergence functions}
\subsubsection{Sum of squared differences}
The sum of the squared differences, denoted as
\begin{equation}    
\Delta_{L^2}(\mathbf{x}, \mathbf{y})=||\mathbf{x}-\mathbf{y}||_2^2=\sum_{i}(x_i-y_i)^2,
\end{equation}
is simple to code, quick to compute and generates similar output to the other divergence function for most test samples, but has little mathematical backing. In particular, it considers all elements of $\mathbf{h}$ to be of equal reliability when, in theory, bins with smaller probabilities should have smaller variances in $\mathbf{h}$.
\subsubsection{Inverse likelihood}
The inverse likelihood function is inversely related to the probability that, under the probabilities produced by the model, the samples $\mathbf{h}$ are drawn. It is given by
\begin{equation}
\Delta_{\mathcal{L}}(\mathbf{x},\mathbf{h})=-\sum_{i}h_i\log(x_i)
\end{equation}
and is derived from the multinomial distribution.
\subsection{Regularizer functions}
\subsubsection{Integral of squared Laplacian}
This is the only regularizer function I have extensively tested. Because the Laplacian is a linear operator, this can be implemented as a matrix operation and an $L^2$ norm.
\begin{equation}
    L_{\nabla^2}(f_{\scriptscriptstyle{D}, \scriptscriptstyle{N}})=\iint\left[\left(\frac{\partial^2}{\partial \delta^2}+\frac{\partial^2}{\partial n^2}\right)f_{\scriptscriptstyle{D}, \scriptscriptstyle{N}}(\delta, n)\right]^2\, dn\,d\delta
\end{equation}
% \subsubsection{Kullback-Leibler divergence}
% The Kullback-Leibler divergence is given by
% \begin{equation}
%     \Delta_{KL}(\mathbf{x},\mathbf{y})=\sum_{i} x_i\log\left(\frac{x_i}{y_i}\right)
% \end{equation}
\printbibliography
\end{document}